\section{Conclusion and Future Work}
\label{sec:conclusion}

In this paper, we addressed the critical challenge of low-data calibration for dynamic model-merging routing heads. We demonstrated that standard classical dynamic routers suffer from severe representational collapse and overfitting on heterogeneous, high-conflict datasets due to sub-optimal optimization in the low-data regime and the zero-sum competitive constraint of the Softmax function.

To resolve these limitations, we evaluated the decoupling of routing pathways via a Softmax-free **Bounded Sigmoidal Router (BSigmoid-Router)** and proposed **Task-Correlation Prior Regularization (TCPR)**, which injects pre-computed task similarity priors (parameter-space or representation-space cosine similarities) into routing head calibration under unit signature normalization and off-diagonal centering.

Through an exceptionally rigorous, seed-controlled empirical evaluation on a heterogeneous four-task Vision Transformer benchmark under computationally constrained expert conditions, we established that:
\begin{itemize}
    \item Decoupling routing pathways via independent sigmoidal projections (\textbf{BSigmoid-Router}) completely eliminates the competitive zero-sum bottleneck of Softmax routing, achieving the top-performing multi-task profile (\textbf{25.50\%} Joint Mean) across all dynamic merging methods.
    \item Under controlled initializations, incorporating static prior regularization (\textbf{TCPR-Param} and \textbf{TCPR-Rep}) fails to deliver empirical improvements over unregularized sigmoidal routing, and actively degrades performance at active scales ($\beta \ge 1.0$).
    \item We deconstruct this failure through the \textbf{Scale Mismatch} (small priors are mathematically dead, rendering $\beta \le 10^{-6}$ identical to the unregularized router) and the \textbf{Alignment-Interference Paradox} (statically forcing weight alignment across disparate domains under sub-optimal expert conditions introduces severe representational noise and interference, collapsing other tasks).
\end{itemize}
These empirical sweeps provide a critical scientific warning to the model-merging community, demonstrating that dynamic calibration is highly robust when unconstrained, and that constraining dynamic routing weights to static, pre-computed priors introduces severe representational damage.

\textbf{Future Work:} In future research, we plan to investigate dynamic, sample-dependent structural regularizations that adaptively align routing pathways on-the-fly during the calibration process to reflect active domain alignment. We also plan to scale sigmoidal dynamic merging to large-scale multi-task benchmarks on modern Large Language Models and visual-language models.
